\begin{lemma}{Eigenschaften der Erweiterung im Vektorraumfall}
             {EigenschaftenDerErweiterungImVektorraumfall}
Seien $K$ ein Körper, $V,W$ $K$--Vektorräume und $n \in \N$.\\
Dann gilt:
\begin{enumerate}[(1)]
	\item \label{EigenschaftenDerErweiterungImVektorraumfall1}
	 Für jedes $\alpha \in \Hom(V, W)$ ist auch 
    $\pfeil \alpha n:V^n \rightarrow W^n$ linear.
\item \label{EigenschaftenDerErweiterungImVektorraumfall2} Die Abbildung \[
\pfeil{\cdot}n: \Hom(V, W) \rightarrow \Hom(V^n, W^n), \varphi\mapsto\pfeil\varphi n\]
ist ebenfalls linear.
\end{enumerate}
\end{lemma}
\begin{proof}
Seien $\varphi, \tau \in \Hom(V, W)$, $k \in K$ und $v, w \in V^n$.\\
Es ist dann
\begin{align*}
\pfeil \varphi n(v+kw)
   & = \pfeil\varphi n \left((v_i+k w_i)_{i=1}^n\right)
     = (\varphi(v_i+kw_i))_{i=1}^n\\
   & = (\varphi(v_i) + k \varphi(w_i))_{i=1}^n &\quad \varphi \text{ linear}\\
   & = \left(\varphi(v_i)\right)_{i=1}^n + k\left(\varphi(v_i)\right)_{i=1}^n
     = \pfeil \varphi n (v) + k \pfeil \varphi n (w)
\end{align*}
Damit ist $\pfeil \varphi n$ linear.\\
Weiter gilt:
\begin{align*}
\pfeil {(\varphi + k \tau)} n (v)
  & = \left((\varphi + k \tau)(v_i)\right)_{i=1}^n
    = \left(\varphi(v_i) + k \tau(v_i)\right)_{i=1}^n\\
  & = \left(\varphi(v_i)\right)_{i=1}^n + k \left(\tau(v_i)\right)_{i=1}^n
    = \pfeil \varphi n (v) + k \pfeil \tau n (v)
\end{align*}
Also ist auch $\pfeil \cdot n$ linear.
\end{proof}